\begin{resumo}

Este trabalho caracteriza a calibração probabilística de um Temporal Fusion Transformer (TFT) aplicado à previsão multi-horizonte de retornos diários, compara-o a uma hierarquia de modelos de referência e descreve a contribuição relativa de famílias de indicadores. O problema de pesquisa decorre do fato de que, para retornos diários de ações líquidas, o componente previsível da média é mínimo, de modo que avaliações restritas a métricas pontuais tendem a medir ruído; por isso, o objeto científico adotado é a distribuição preditiva, e não a acurácia do ponto previsto. O recorte é \textit{single-asset} (um modelo por ativo), com AAPL como ativo piloto e horizontes $h+1$ e $h+7$. A metodologia integra dados de mercado, indicadores técnicos, sentimento de notícias estimado por FinBERT e fundamentos contábeis, com causalidade temporal verificada \textit{feature} a \textit{feature}, e treina o TFT em modo quantílico sobre uma grade densa de sete níveis, com rearranjo monotônico dos quantis. O protocolo experimental usa validação \textit{walk-forward} de janela expansiva com purga e embargo em dias de pregão e conjunto de calibração dedicado, seis \textit{folds} e dez sementes, alinhamento estrito por \textit{target timestamp} e comparação contra uma hierarquia pré-declarada de \textit{baselines}: retorno nulo, média histórica, AR(1), EWMA-vol, quantis empíricos históricos e \textit{gradient boosting} quantílico. A calibração é julgada, por horizonte, por um \textit{gate} de bandas de Wilson nas duas caudas do intervalo central de 80\%, sempre acompanhado do poder declarado contra desvios mínimos pré-registrados. A \textit{pinball loss} é a métrica primária e o CRPS, complementar; a inferência usa o teste de Diebold-Mariano (HAC/HLN, unilateral) com ajuste de Holm e o \textit{Model Confidence Set}, e a calibração nativa é confrontada com \textit{conformalized quantile regression} como referência de cobertura empírica. A contribuição de famílias é lida pela concordância entre importância por permutação e ablação com re-treino, com os pesos da \textit{Variable Selection Network} como descrição complementar. O veredito é emitido por um \textit{scorecard} aplicado mecanicamente sob um pré-registro congelado, \textit{hasheado} e ancorado antes da geração das tabelas que o veredito julga, e a refutação das hipóteses é tratada como resultado válido. Na rodada confirmatória, a calibração do intervalo central de 80\% não foi rejeitada em nenhum dos dois horizontes, com poder superior a \(0{,}97\) contra desvios de 5 pontos percentuais na cobertura; as violações, contudo, agrupam-se no tempo, de modo que a calibração é marginal, e não condicional. O TFT superou os \textit{baselines} naive nos dois horizontes, mas não os comparadores fortes: em \(h+1\) foi excluído do \textit{Model Confidence Set}, e em \(h+7\) empatou com os fortes, sendo declarado vencedor primário pela regra pré-registrada, embora a menor perda média tenha sido a do EWMA-vol nos dois horizontes. A contribuição das famílias não variou de forma detectável entre horizontes, e a família técnica, que reúne as medidas de volatilidade, foi a mais importante nos dois métodos. Os resultados indicam que o valor do modelo está na descrição calibrada da escala da distribuição, e não em ganho de \textit{skill} sobre modelos de volatilidade simples.

\textbf{[Nota de versão — a ser removida na versão final de defesa]} Este resumo é uma versão em evolução, atualizada a cada rodada de orientação conforme a consolidação dos resultados confirmatórios. A versão final, sem esta nota, será aquela submetida à banca avaliadora.

\textbf{Palavras-chave}: Temporal Fusion Transformer. Previsão quantílica. Calibração probabilística. Predição conformal. Avaliação fora da amostra.

\end{resumo}
